\scanpage{9}

% archive sheet number: 8

\uline{Supposons que} $\lambda\to\lambda1_C$ soit un isomorphisme de $A$ sur un facteur direct du $A$-module $C$ [pour ceci, il suffit, comme on voit facilement, que $C$ soit un $A$-module \deleted{libre de type} \inserted{ayant} une base contenant $1_C$, p.ex. que ce soit un $A$-module libre de type fini]. Alors la condition précédente implique facilement $x=\lambda1_C$, où $\lambda\in A$ est déterminé de façon unique. Dans ce cas, on a donc bien ce qu'on voulait.

\uline{N.B.} En général, \deleted{\unsure{on}} il ne sera pas vrai que le quotient de $C$ par $B$ (ds $\unsure{\mathcal{C}|\mathscr{C}}$) soit $A$, sous les conditions indiquées. On peut prendre par exemple $C=B/\unsure{J|I}B$, où $J$ est un idéal non nul de $A$, \deleted{\unsure{Alors}} \deleted{de définir $q_C$} en prenant la structure d'algèbre quotient, et définissant $q_C:C\to B\otimes C$ i.e.

\[
q_C:B/JB\longrightarrow (B\otimes B)/J.(B\otimes B)
\]

en réduisant $\unsure{p|\rho}_B$ mod $J$. Les axiomes restent vérifiés, cependant l'algèbre \unsure{semi-primitive|des primitifs} de $C$ est $A/J$, et non $A$\ldots\ldots
